\section{Introduction}
\label{sec:introduction}

A binary matrix has both a row order and a column order. For a fixed
binary $k\times k$ matrix $H$, an \emph{ordered copy} of $H$ in a binary
$n\times n$ matrix $A$ is a choice
\[
 r_1<\cdots<r_k,\qquad c_1<\cdots<c_k
\]
with $A(r_i,c_j)=H(i,j)$ for every $i,j\in[k]$, where
$[k]=\{1,\ldots,k\}$. Thus a copy matches the zeros as well as the ones,
and the row and column choices are independent. Write $N_H(A)$ for the
number of such copies and call $A$ \emph{$H$-free} if $N_H(A)=0$.
The normalized distance from $H$-freeness is
\[
 \operatorname{dist}_H(A)
 =\frac1{n^2}\min_{B\text{ is }H\text{-free}}
   |\{(r,c)\in[n]^2:A(r,c)\ne B(r,c)\}|,
\]
where the minimum is over binary $n\times n$ matrices. Every cell may be
changed, in either direction. We say that $A$ is $\epsilon$-far from
being $H$-free when $\operatorname{dist}_H(A)\ge\epsilon$.

The polynomial ordered binary matrix-removal conjecture asserts that
for each fixed $H$ there are constants $c_H,C_H>0$ such that, for every
integer $n\ge1$, every $0<\epsilon<1$, and every binary $n\times n$ matrix $A$,
\[
 N_H(A)\ge c_H\epsilon^{C_H}n^{2k}
 \quad\text{whenever}\quad \operatorname{dist}_H(A)\ge\epsilon.
\]
We disprove this conjecture with an explicit pattern of size $66$.

\begin{theorem}\label{thm:main}
There is a fixed binary $66\times66$ matrix $H$ with the following
property. For every integer $h\ge1$, put
\[
 n_h=(386h+2)2^h,\qquad \epsilon_h=(386h+2)^{-2}.
\]
There is a binary $n_h\times n_h$ matrix $A_h$ such that
\[
 \operatorname{dist}_H(A_h)\ge\epsilon_h,
 \qquad
 \frac{N_H(A_h)}{n_h^{132}}\le\epsilon_h2^{-h}.
\]
Consequently, no positive constants $c,C$ give a polynomial removal
bound for this fixed $H$.
\end{theorem}

The examples also give a lower bound for canonical sampling
(Corollary~\ref{cor:canonical-testing}). A tester that independently
chooses uniform $q$-element row and column sets and rejects only when
the sampled ordered submatrix contains $H$ needs
$q\ge\exp(\Omega(\epsilon_h^{-1/2}))$ for any fixed positive rejection
probability along this sequence. Here $q$ is the number of sampled
positions on each axis; the submatrix has $q^2$ entries.

\paragraph{History and significance.}
Alon, Fischer and Newman raised the ordered-matrix question after proving
polynomial-query testing for fixed finite forbidden families when row
and column order is ignored~\cite[Section~7]{AlonFischerNewman2007}.
Alon and Ben-Eliezer formulated the polynomial removal question for
finite families of binary matrices~\cite[Problem~1.4]{AlonBenEliezer2020}.
The singleton version appears explicitly as
Conjecture~4.5 in the first arXiv version of the survey by Gishboliner and
Shapira~\cite{GishbolinerShapira2025}.
Theorem~\ref{thm:main} resolves that conjecture negatively and answers
the finite-family question negatively already for one square pattern.

Qualitative ordered matrix removal was established by Alon,
Ben-Eliezer and Fischer~\cite[Theorem~1.8 in the full
version]{AlonBenEliezerFischer2017}: a positive copy-density bound exists
for each fixed positive distance. Our examples concern its quantitative
dependence. Polynomial bounds also hold for finite families closed under
row permutations, or under column permutations by
transposition~\cite[Theorem~1.5]{AlonBenEliezer2020}.
The polynomial induced removal classification for ordered graphs of Gishboliner
and Tomon~\cite[Theorem~1.1]{GishbolinerTomon2022} concerns a
different model, with a single vertex order and symmetric adjacency.

Nonpolynomial lower bounds were already known for ternary hosts:
Fischer and Rozenberg obtained them for a fixed two-symbol $2\times2$
pattern, with row and column order ignored~\cite{FischerRozenberg2007}.
Alon and Ben-Eliezer give a related
ternary-matrix construction based on progression-free
sets~\cite[Section~5.1]{AlonBenEliezer2020}. These examples use a third
symbol in the host matrix; our construction uses only $0$ and $1$.

A particularly relevant positive result is the polynomial lower bound
for matrices containing many entry-disjoint copies of a fixed
pattern~\cite[Theorem~1.2]{AlonBenEliezer2020}. It yields a small set of
cells meeting all existing copies when the copy density is small.
Changing those cells can create new copies. Our construction makes
this distinction explicit: all original copies meet a set of $2^h$
cells, while any binary repair changes at least $4^h$ cells.

\paragraph{Proof strategy.}
The last two rows and columns of $H$ meet in the pattern
\[
 P=\begin{pmatrix}1&0\\1&1\end{pmatrix}.
\]
Its first $64$ rows and columns specify which pairs of rows and columns
can play these two roles. The host matrix contains many copies of this
$64\times64$ anchor, each imposing a different instance of the same
$2\times2$ restriction. Section~\ref{sec:construction} gives the pattern
and the host entries explicitly.

The host has rows and columns indexed by a binary tree, with a plus
block and a minus block at each internal node on each axis. At each
leaf, the two blocks on each axis coincide. One root entry is $1$ and
the other is $0$. Avoiding $P$ in the designated bodies propagates
these values down a common path, first to the next row level and then
to the next column level. At the leaves the two signs share their
positions, so they cannot retain opposite values.

We first prove the distance bound in Section~\ref{sec:distance}.
Fix an arbitrary edited matrix and sample a tree path together with
anchor positions and auxiliary ``dummy'' positions, whose entries
against the tree positions are all $1$ in the original host. Each axis is
partitioned into
classes of equal size $m=2^h$. For every cell needed to keep the sampled
anchors, dummy incidences and root values unchanged, its row and column
are selected with probability $m^{-2}$. Summing over the actual edited
cells shows that fewer than $m^2$ changes leave a suitable selection.
All other entries may have changed; avoidance of $H$ itself forces the
contradictory propagation in the edited matrix. This cellwise sampling
argument avoids a loss depending on the number of classes and gives
the inverse-square distance in Theorem~\ref{thm:main}.

The copy bound concerns the unedited host and requires locating every
copy, beyond the designated anchors used in the distance proof.
In Section~\ref{sec:copies}, dense rows and columns, together with a
five-column block whose rows realize all binary words, force every
anchor copy to occupy one of the prescribed anchor blocks. The remaining rows and columns must then have
that block's designated roles. The order and threshold rules then force
every copy to use a diagonal leaf cell, giving the small copy count.
Section~\ref{sec:conclusion} combines the two estimates and derives the
sampling consequence.
